% #############################################################################
% RESUMO em Português
% !TEX root = ../main.tex
% #############################################################################
% reset acronyms
\acresetall
% O guia limita o resumo a 250 palavras.
% use \noindent in firts paragraph
\noindent A Inteligência Artificial Generativa e os Modelos de Linguagem de Grande Dimensão passaram de assistentes de preenchimento automático a contribuidores activos no Ciclo de Vida de Desenvolvimento de Software, permitindo exprimir requisitos em linguagem natural e supervisionar os artefactos gerados. A adopção, contudo, ultrapassou o método. Uma Revisão Sistemática da Literatura sobre 646 artigos confirma benefícios substanciais a par de riscos recorrentes de fiabilidade, segurança, rastreabilidade e responsabilização, e identifica a ausência de orientação estruturada e repetível para integrar estas ferramentas em contextos profissionais como o principal problema por resolver. A maioria do código gerado por IA funcionalmente correcto contém defeitos de segurança, e os estrangulamentos deslocam-se para a revisão em vez de desaparecerem: o valor depende do processo, não do modelo.

Esta dissertação aborda essa lacuna através da Design Science Research Methodology, propondo uma framework orientada a processo para a colaboração Humano-IA governada, cujas contribuições de IA devem ser delimitadas, validadas e rastreadas. Define fases com portões de qualidade específicos, responsabilidades expressas como funções e não como cargos, e métricas de governação, ancoradas numa especificação versionada que liga cada artefacto ao seu requisito. É instanciada no Apex, aplicação web que operacionaliza cada fase, portão e artefacto sobre uma ferramenta real de gestão de projectos. A avaliação combina aplicação num projecto real, instrumentos padronizados de usabilidade e carga cognitiva, entrevistas, e avaliação analítica face a critérios previamente definidos. O contributo é um modelo de processo fundamentado em evidência que permite adoptar assistência por IA sem abdicar do controlo sobre qualidade, segurança e responsabilização.
